% ============================================================
% Report — Domain Adaptability in Video Anomaly Detection:
%          Literature Landscape, Model Choices, and Gap Analysis
% Prepared in response to the guide's directive (July 2026):
%   (1) research the domain-adaptability literature more deeply,
%   (2) identify which models each line of work uses,
%   (3) identify gaps, and justify which model we will use.
% ============================================================
\documentclass[12pt, a4paper]{article}

% --- Packages ---
\usepackage[utf8]{inputenc}
\usepackage[T1]{fontenc}
\usepackage{lmodern}
\usepackage[margin=1in]{geometry}
\usepackage{hyperref}
\usepackage{xcolor}
\usepackage{enumitem}
\usepackage{booktabs}
\usepackage{titlesec}
\usepackage{fancyhdr}
\usepackage{tabularx}
\usepackage{array}
\usepackage{parskip}
\usepackage{longtable}

% --- Colors ---
\definecolor{sectionblue}{HTML}{1A5276}
\definecolor{linkblue}{HTML}{2471A3}
\definecolor{highlightbg}{HTML}{EBF5FB}
\definecolor{accentgreen}{HTML}{1E8449}

% --- Hyperlink Style ---
\hypersetup{
    colorlinks=true,
    linkcolor=linkblue,
    citecolor=linkblue,
    urlcolor=linkblue,
    pdftitle={Domain Adaptability in Video Anomaly Detection},
    pdfauthor={Prakash Kumar Sarangi}
}

% --- Section Formatting ---
\titleformat{\section}{\Large\bfseries\color{sectionblue}}{\thesection}{0.6em}{}[\titlerule]
\titleformat{\subsection}{\large\bfseries\color{sectionblue}}{\thesubsection}{0.6em}{}
\titleformat{\subsubsection}{\normalsize\bfseries}{\thesubsubsection}{0.6em}{}

% --- Header / Footer ---
\pagestyle{fancy}
\fancyhf{}
\lhead{\small\textcolor{gray}{Domain Adaptability in VAD --- Findings Report}}
\rhead{\small\textcolor{gray}{\thepage}}
\renewcommand{\headrulewidth}{0.4pt}

\newcolumntype{L}[1]{>{\raggedright\arraybackslash}p{#1}}

% ============================================================
\begin{document}

% ============ TITLE PAGE ============
\begin{titlepage}
\centering
\vspace*{1.5cm}

{\large National Institute of Technology Calicut}\\[0.2cm]
{\normalsize Department of Computer Science and Engineering}\\[1.8cm]

{\Large\bfseries Domain Adaptability in Video Anomaly Detection:}\\[0.4cm]
{\Large\bfseries Literature Landscape, Model Choices, and Gap Analysis}\\[1.5cm]

{\normalsize\itshape
A report prepared in response to the guide's directive to research the
domain-adaptability literature in greater depth, identify the models used
by existing approaches, identify open gaps, and justify the model chosen
for this thesis.}\\[2cm]

\begin{tabular}{rl}
    \textbf{Student:}     & Prakash Kumar Sarangi \\
    \textbf{Programme:}   & M.Tech, Computer Science and Engineering \\
    \textbf{Roll No.:}    & \underline{\hspace{4cm}} \\
    \textbf{Guide:}       & \underline{\hspace{4cm}} \\
    \textbf{Date:}        & \today \\
\end{tabular}

\vfill
{\small Working title of the thesis project: \textbf{DA-ZVAD} --- Domain-Adaptive
Zero-shot Video Anomaly Detection}
\end{titlepage}

% ============================================================
\section{Objective of This Report}
% ============================================================

Following the discussion with the guide, this report addresses three specific
requests: (i) a deeper survey of the literature on \textbf{domain adaptability}
in video anomaly detection (VAD), (ii) an explicit account of \textbf{which
models} each line of work relies upon, and (iii) an identification of
\textbf{gaps} in the literature together with a justification of the model
combination chosen for this thesis, DA-ZVAD.

\textbf{Scope of ``domain adaptability.''} In this literature, the term covers
a spectrum of approaches --- from fully supervised adversarial domain
adaptation, through few-shot / meta-learned scene adaptation, to fully
training-free transfer via frozen vision-language models (VLMs) and large
language models (LLMs). This report organises the literature along that
spectrum rather than treating it as a single homogeneous category, since the
cost, assumptions, and applicability of each family differ substantially.

% ============================================================
\section{The Domain-Adaptability Landscape: Four Lanes}
% ============================================================

The reviewed literature falls into four lanes, ordered by decreasing cost of
adapting to a new (target) domain.

\begin{table}[h!]
\centering
\renewcommand{\arraystretch}{1.35}
\begin{tabularx}{\textwidth}{@{}l L{3.1cm} X@{}}
\toprule
\textbf{Lane} & \textbf{Target-domain cost} & \textbf{Representative work} \\
\midrule
A. Learned / adversarial DA & Target data + a training run & Ada-VAD, DA-VAD, Industrial-AdaVAD \\
B. Meta-learning / few-shot scene adaptation & A few target frames + gradient steps & Lu \textit{et al.} (ECCV'20), Anomaly Crossing, adversarial-diffusion few-shot VAD \\
C. Zero-shot cross-domain (vision-only) & None at target time (source training still required) & zxVAD (WACV'23) \\
D. Training-free VLM/LLM-based & None --- all components frozen at inference & LAVAD, VERA, AnyAnomaly, MM-VAD (2026) \\
\bottomrule
\end{tabularx}
\caption{The four lanes of domain-adaptability approaches in VAD, by adaptation cost.}
\end{table}

\subsection{Lane A --- Learned / Adversarial Domain Adaptation}

\textbf{Ada-VAD} (SIAM SDM 2024) tackles the few-shot cross-domain VAD problem
in two stages: a pretraining stage that synthesises abnormal samples and
trains a domain-invariant predictor using a self-supervised prediction task,
followed by an adaptation stage that fine-tunes on a handful of target-domain
frames using adversarial training to mitigate distribution shift.
\textit{Model used:} a 3D-convolutional / autoencoder-style future-frame
predictor --- no vision-language model is involved.

\textbf{DA-VAD} (2025) is weakly supervised and treats video-level and
clip-level features as distinct source and target domains, inserting explicit
\textit{domain aligner} modules between them to reduce overfitting to a
specific scene's visual statistics. \textit{Model used:} I3D-style clip
features with alignment heads.

\textbf{Industrial-AdaVAD} (\textit{Mathematics}, 2025) targets
resource-constrained industrial deployments and proposes a multilayer
adversarial domain-adaptation mechanism for cross-scene transfer at the
edge. \textit{Model used:} a lightweight CNN backbone with adversarial
discriminators.

\textit{Common limitation of Lane A:} each new deployment requires target
data and an explicit training run; the adapted model becomes scene-specific;
none of these methods produces natural-language explanations of detected
anomalies.

\subsection{Lane B --- Meta-Learning and Few-Shot Scene Adaptation}

The foundational work here is \textbf{Lu \textit{et al.}, ``Few-shot
Scene-adaptive Anomaly Detection''} (ECCV 2020), which applies MAML-style
meta-learning across many training scenes so that a future-frame-prediction
GAN can adapt to an unseen scene from only a handful of frames.
\textbf{Anomaly Crossing} explicitly reframes cross-domain VAD as a
cross-domain few-shot learning problem, fusing source and target
representations. A more recent \textbf{adversarial-diffusion} approach
(\textit{Neurocomputing}, 2024) replaces the GAN generator with a diffusion
model adapted per scene from few samples.

\textit{Common limitation of Lane B:} adaptation still requires gradient
updates per scene, and the backbone models predate the VLM era (GANs,
autoencoders, 3D-CNNs) --- there is no language interface and no
explanation capability.

\subsection{Lane C --- Zero-Shot Cross-Domain, Vision-Only}

\textbf{zxVAD} (WACV 2023) is the closest \textit{non-VLM} relative of the
present thesis: it performs cross-domain VAD \textbf{without any target-domain
adaptation}. A \textit{Normalcy Classifier} learns to contrast normal-event
features against pseudo-abnormal features synthesised by an
\textit{untrained} convolutional network that pastes foreign objects into
normal frames --- avoiding the cost of collecting real anomalies. The
backbone remains a future-frame-prediction network.

\textit{Limitation:} the source model is still trained (on the source
domain); ``normal'' is defined purely by learned visual statistics, with no
mechanism to \textit{tell} the system what normal means in a new scene; no
explanations are produced; and evaluation stays within the surveillance
domain family (no industrial-to-surveillance transfer is studied).

\subsection{Lane D --- Training-Free VLM/LLM Methods (2024--2026): The Active Frontier}

This is the most active lane and the one this thesis builds on directly.

\textbf{LAVAD} (CVPR 2024) is fully training-free: a BLIP-2-based captioner
converts frames to text, a Llama-family LLM scores temporal windows of these
captions for anomalousness, and an ImageBind-style cross-modal similarity
step cleans noisy captions. It is strong on UCF-Crime and XD-Violence but
chains three large frozen models and has no explicit mechanism for injecting
domain context.

\textbf{VERA} (CVPR 2025) keeps its VLM backbone (InternVL2) entirely frozen
and instead makes only a set of \textit{guiding questions} learnable; these
questions are optimised offline, on source data, through a learner/optimiser
VLM interaction, and the model is explainable by construction. A key
observation for our positioning: VERA's verbalisation is \textbf{learned at
training time on the source domain}, not adapted at test time on a new
domain --- cross-domain transfer of the learned questions is not the problem
VERA studies.

\textbf{AnyAnomaly} (2025) introduces \textit{customizable} VAD: a
user-supplied text description of what should count as an anomaly drives a
zero-shot large vision-language model (LVLM) through segment-level,
context-aware visual question answering, with no training at all. It is the
clearest precedent for using text to define the task at inference time,
though it verbalises the \textit{anomaly}, not the \textit{scene/normality}.

\textbf{MM-VAD} (2026) reframes training-free VAD as \textit{adaptive
test-time inference}: caption-derived scene representations are projected
into hyperbolic space and reasoned over by a frozen LLM via adaptive
question answering. Together with several concurrent 2025--2026 works ---
MoniTor (online training-free VAD), MLLM-based zero-shot VAD without real
anomaly examples, sparse-reasoning frameworks over frozen large models, and
latent-manifold steering in frozen MLLMs --- this confirms that the field is
moving decisively toward training-free, test-time reasoning pipelines.
Weakly supervised CLIP-prompt approaches such as VadCLIP and HeadCLIP sit
between Lanes A and D: CLIP-based, but with learnable prompt tokens that
still require training.

\textit{Common blind spot of Lane D:} every method in this lane is developed
and evaluated \textbf{within the surveillance-video domain only}. None
isolates domain adaptation as the measured variable --- there is no study
that (i) moves a single training-free pipeline across genuinely different
domain families (industrial inspection and surveillance video), (ii) injects
the target domain purely as a \textbf{test-time text description} rather
than a learned prompt or question set, and (iii) ablates \textbf{which
component of the pipeline actually carries the cross-domain transfer}.

% ============================================================
\section{Which Models Each Line of Work Uses}
% ============================================================

Table~\ref{tab:models} summarises, for every reviewed line of work, the
backbone model(s) used, whether any component is trained, whether a language
channel exists for injecting domain context, and whether the method produces
explanations. The final row states the model combination used in this
thesis (DA-ZVAD) for direct comparison.

\begin{table}[h!]
\centering
\renewcommand{\arraystretch}{1.4}
\small
\begin{tabularx}{\textwidth}{@{}L{3.4cm} X c c c@{}}
\toprule
\textbf{Work} & \textbf{Backbone(s)} & \textbf{Trained?} & \textbf{Language channel} & \textbf{Explains?} \\
\midrule
Ada-VAD / DA-VAD / Industrial-AdaVAD & 3D-CNN / I3D / autoencoder + adversarial heads & Yes (source \& target) & None & No \\
Lu \textit{et al.} (ECCV'20) & Future-frame-prediction GAN + MAML & Yes (meta + per-scene) & None & No \\
zxVAD & Frame-prediction CNN + untrained-CNN anomaly synthesis & Yes (source only) & None & No \\
LAVAD & BLIP-2 + Llama-family LLM + ImageBind & No & Captions (implicit) & Partial \\
VERA & InternVL2 (frozen); guiding questions learned & Verbal parameters only & Learned guiding questions & Yes \\
AnyAnomaly & Open LVLM with VQA & No & User-defined anomaly text & Partial \\
MM-VAD (2026) & Frame captioner + frozen LLM (hyperbolic reps.) & No & Adaptive QA & Partial \\
\midrule
\textbf{DA-ZVAD (this thesis)} & \textbf{Frozen OpenCLIP ViT-L/14 + LLaVA-1.5-7B (4-bit)} & \textbf{No --- nothing, anywhere} & \textbf{Test-time verbalised scene context} & \textbf{Yes} \\
\bottomrule
\end{tabularx}
\caption{Model choices across the reviewed literature, compared with the model combination used in this thesis.}
\label{tab:models}
\end{table}

% ============================================================
\section{Identified Gaps}
% ============================================================

Five gaps emerge from the comparison above.

\begin{enumerate}[leftmargin=1.6em, itemsep=0.5em]
    \item \textbf{G1 --- No cross-domain-family evaluation.} No VLM-based
    method is evaluated across genuinely different domain families
    (industrial inspection \textit{and} surveillance video) with a single
    frozen pipeline. zxVAD achieves zero-shot transfer but stays within
    surveillance-style domains; every Lane D method is surveillance-only.

    \item \textbf{G2 --- No test-time, language-only domain injection.} VERA
    verbalises \textit{scene reasoning} but the verbalisation is learned
    offline on source data; AnyAnomaly verbalises the \textit{anomaly}, not
    the \textit{scene}. No method describes the target scene/normality in
    text purely at test time, with zero optimisation, and measures the
    resulting effect under domain shift.

    \item \textbf{G3 --- No component-wise transfer ablation.} No existing
    study quantifies \textit{which} component of a training-free pipeline
    (visual backbone, prompt wording, verbalised context, or temporal
    aggregation) is responsible for --- or limits --- cross-domain transfer.

    \item \textbf{G4 --- Explanation quality under domain shift is
    unstudied.} VERA demonstrates explainability, but only within its
    (single) training domain; how explanation quality behaves when the
    scene changes has not been examined anywhere in the literature reviewed.

    \item \textbf{G5 --- No minimal, single-GPU, fully training-free
    stack.} LAVAD, the strongest fully training-free baseline, chains three
    large models (captioner, LLM, cross-modal refiner). A minimal
    training-free stack that still supports domain adaptation and
    explanation, and fits a single consumer GPU, does not yet exist in the
    literature.
\end{enumerate}

% ============================================================
\section{Model Justification: Why CLIP ViT-L/14 + LLaVA-1.5-7B for DA-ZVAD}
% ============================================================

DA-ZVAD is designed to target G1--G5 jointly, using a single frozen pipeline:
a frozen OpenCLIP ViT-L/14 visual encoder for zero-shot anomaly scoring, a
training-free temporal aggregation stage, a \textbf{test-time verbalised
scene-context} module (in contrast to VERA's offline-learned questions), and
a frozen LLaVA-1.5-7B (4-bit) reasoning stage for natural-language
explanations. Every stage remains frozen; nothing is trained anywhere in the
pipeline.

The specific model choice is justified as follows:

\begin{itemize}[leftmargin=1.6em, itemsep=0.4em]
    \item \textbf{CLIP ViT-L/14} is the common visual backbone of the
    zero-shot anomaly-detection literature (the WinCLIP lineage, AnomalyCLIP,
    and related work), which keeps our cross-domain AUROC numbers directly
    comparable to published baselines rather than confounded by a different
    encoder.
    \item \textbf{LLaVA-1.5-7B in 4-bit quantisation} is the lightest capable
    open multimodal LLM available (approximately 5.5\,GB of VRAM), which lets
    the entire stack --- encoder, temporal module, context module, and
    reasoning module --- run on a single consumer-grade GPU (a T4-class or
    RTX-3060-class card), in deliberate contrast to LAVAD's three-model
    chain (addressing G5).
    \item \textbf{Keeping every component frozen} makes the domain-adaptation
    claim clean and attributable: since no parameters are updated anywhere,
    any measured change in performance when the verbalised scene description
    changes can be attributed to the \textit{text}, not to hidden weight
    updates --- directly addressing G2 and enabling the component ablation
    of G3.
\end{itemize}

% ============================================================
\section{Next Steps}
% ============================================================

\begin{enumerate}[leftmargin=1.6em, itemsep=0.4em]
    \item Run the zero-shot detection probe on real ShanghaiTech video to
    obtain the first frame-level AUROC number for the training-free
    pipeline (Sem-3 Phase~A deliverable).
    \item Execute the industrial\,$\leftrightarrow$\,surveillance
    cross-domain evaluation (MVTec AD $\leftrightarrow$ ShanghaiTech / CUHK
    Avenue) to directly address G1.
    \item Run the module on/off ablation grid (visual backbone, temporal
    aggregation, verbalised context, reasoning) to address G3.
    \item Evaluate explanation quality specifically under domain shift
    (source-domain vs.\ target-domain prompts) to address G4.
    \item Incorporate guide feedback on this report and finalise the
    positioning of the cross-domain, test-time verbalised-context
    contribution (G2) as the central claim of the thesis.
\end{enumerate}

% ============================================================
\section{References}
% ============================================================

\begin{enumerate}[label={[\arabic*]}, leftmargin=2.2em, itemsep=0.3em]

\item S.\ Aich, J.-C.\ Chiang, and P.\ Wang, ``Cross-Domain Video Anomaly
Detection without Target Domain Adaptation,'' in \textit{Proc.\ IEEE/CVF
WACV}, 2023. \href{https://arxiv.org/abs/2212.07010}{arXiv:2212.07010}

\item ``Ada-VAD: Domain Adaptable Video Anomaly Detection,'' in
\textit{Proc.\ SIAM Int.\ Conf.\ on Data Mining (SDM)}, 2024.
\href{https://epubs.siam.org/doi/10.1137/1.9781611978032.73}{doi:10.1137/1.9781611978032.73}

\item Y.\ Lu, F.\ Yu, M.\ K.\ K.\ Reddy, and Y.\ Wang, ``Few-Shot
Scene-Adaptive Anomaly Detection,'' in \textit{Proc.\ ECCV}, 2020.

\item ``Anomaly Crossing: New Horizons for Video Anomaly Detection as
Cross-Domain Few-Shot Learning,'' 2021.
\href{https://arxiv.org/abs/2112.06320}{arXiv:2112.06320}

\item ``Adversarial Diffusion for Few-Shot Scene-Adaptive Video Anomaly
Detection,'' \textit{Neurocomputing}, 2024.

\item ``Industrial-AdaVAD: Adaptive Industrial Video Anomaly Detection
Empowered by Edge Intelligence,'' \textit{Mathematics}, vol.\ 13, no.\ 17,
2025.

\item L.\ Zanella, W.\ Menapace, M.\ Mancini, Y.\ Wang, and E.\ Ricci,
``Harnessing Large Language Models for Training-Free Video Anomaly
Detection,'' in \textit{Proc.\ IEEE/CVF CVPR}, 2024.
\href{https://arxiv.org/abs/2404.01014}{arXiv:2404.01014}

\item M.\ Ye, W.\ Liu, and P.\ He, ``VERA: Explainable Video Anomaly
Detection via Verbalized Learning of Vision-Language Models,'' in
\textit{Proc.\ IEEE/CVF CVPR}, 2025.
\href{https://arxiv.org/abs/2412.01095}{arXiv:2412.01095}

\item ``AnyAnomaly: Zero-Shot Customizable Video Anomaly Detection with
LVLM,'' 2025. \href{https://arxiv.org/abs/2503.04504}{arXiv:2503.04504}

\item ``MM-VAD: Geometry-Aware Semantic Reasoning for Training-Free Video
Anomaly Detection,'' 2026.
\href{https://arxiv.org/abs/2603.13374}{arXiv:2603.13374}

\item X.\ Ding, ``Quo Vadis, Anomaly Detection? LLMs and VLMs in the
Spotlight,'' 2024.
\href{https://arxiv.org/abs/2412.18298}{arXiv:2412.18298}

\item P.\ Wu \textit{et al.}, ``Open-Vocabulary Video Anomaly Detection,''
in \textit{Proc.\ IEEE/CVF CVPR}, 2024.

\end{enumerate}

\end{document}
